\documentclass[11pt,letterpaper]{article}

% Packages
\usepackage[utf8]{inputenc}
\usepackage[margin=1in]{geometry}
\usepackage{graphicx}
\usepackage{hyperref}
\usepackage{listings}
\usepackage{xcolor}
\usepackage{fancyhdr}
\usepackage{titlesec}
\usepackage{tocloft}
\usepackage{tabularx}
\usepackage{longtable}
\usepackage{booktabs}
\usepackage{enumitem}
\usepackage{microtype}
\usepackage{float}

% Hyperref setup
\hypersetup{
    colorlinks=true,
    linkcolor=blue,
    filecolor=magenta,
    urlcolor=cyan,
    pdftitle={Blockd Platform Architecture},
    pdfauthor={Blockd Development Team},
}

% Code listing setup
\definecolor{codegreen}{rgb}{0,0.6,0}
\definecolor{codegray}{rgb}{0.5,0.5,0.5}
\definecolor{codepurple}{rgb}{0.58,0,0.82}
\definecolor{backcolour}{rgb}{0.95,0.95,0.92}

\lstdefinestyle{mystyle}{
    backgroundcolor=\color{backcolour},
    commentstyle=\color{codegreen},
    keywordstyle=\color{magenta},
    numberstyle=\tiny\color{codegray},
    stringstyle=\color{codepurple},
    basicstyle=\ttfamily\footnotesize,
    breakatwhitespace=false,
    breaklines=true,
    captionpos=b,
    keepspaces=true,
    numbers=left,
    numbersep=5pt,
    showspaces=false,
    showstringspaces=false,
    showtabs=false,
    tabsize=2
}
\lstset{style=mystyle}

% Header and footer
\pagestyle{fancy}
\fancyhf{}
\rhead{Blockd Platform Architecture}
\lhead{Version 1.0.0}
\cfoot{\thepage}

% Title page
\title{
    \Huge\textbf{Blockd Interview Security Platform}\\
    \vspace{0.5cm}
    \Large Complete System Architecture \& Deployment Guide\\
    \vspace{1cm}
    \includegraphics[width=0.3\textwidth]{placeholder_logo}\\[1cm]
}
\author{Blockd Development Team}
\date{November 2025\\Version 1.0.0}

\begin{document}

\maketitle
\newpage

% Table of Contents
\tableofcontents
\newpage

% List of Tables and Figures
\listoftables
\listoffigures
\newpage

%----------------------------------------------------------------------------------------
% EXECUTIVE SUMMARY
%----------------------------------------------------------------------------------------

\section{Executive Summary}

\subsection{Project Overview}

Blockd is an enterprise-grade interview security and AI-powered anti-cheating platform designed to ensure the integrity of technical interviews conducted remotely. The platform combines advanced security monitoring, artificial intelligence-based answer detection, real-time eye tracking, and comprehensive session management to provide a robust solution for organizations conducting thousands of concurrent interviews.

\subsection{Key Capabilities}

\begin{itemize}[noitemsep]
    \item \textbf{Custom Chromium Browser}: Native browser fork with embedded security monitoring that cannot be bypassed
    \item \textbf{AI-Powered Detection}: Multi-LLM comparison (GPT-4, Claude, Gemini) with XGBoost ensemble classification
    \item \textbf{Eye Tracking}: Real-time gaze analysis at 30 FPS with MediaPipe FaceMesh and off-screen detection
    \item \textbf{Security Monitoring}: Process detection, VM detection, window tracking, and clipboard monitoring
    \item \textbf{Real-time Video}: WebRTC streaming with mediasoup SFU supporting 100+ concurrent streams
    \item \textbf{Comprehensive Reporting}: Automated risk scoring, detailed analytics, and PDF report generation
\end{itemize}

\subsection{Technology Highlights}

\begin{table}[H]
\centering
\begin{tabularx}{\textwidth}{lX}
\toprule
\textbf{Category} & \textbf{Technologies} \\
\midrule
Frontend & React 19.2.0, TypeScript 5.9.3, Vite 6.x, Tailwind CSS 4.x \\
Backend & Node.js 24.11.0, Python 3.14.0, Fastify 5.x, FastAPI 0.121.3 \\
Databases & PostgreSQL 18.1, Redis 8.4, TimescaleDB 2.x, RabbitMQ 4.x \\
ML/AI & PyTorch 2.5.x, MediaPipe 0.10.x, OpenAI GPT-4, Claude 3.5, Gemini 1.5 \\
Browser & Chromium 142.0.7444.175 (custom fork) \\
Infrastructure & Docker 27.x, Kubernetes 1.31, Terraform, Helm 3.x \\
\bottomrule
\end{tabularx}
\caption{Core Technology Stack}
\end{table}

\subsection{Implementation Metrics}

\begin{itemize}[noitemsep]
    \item \textbf{Total Files}: 674+ production files
    \item \textbf{Lines of Code}: 91,000+ lines
    \item \textbf{Microservices}: 8 backend services fully implemented
    \item \textbf{Development Phases}: 6 phases
    \item \textbf{Test Coverage}: 68+ test scenarios, E2E, integration, and load tests
    \item \textbf{Documentation}: 16+ comprehensive specification documents
\end{itemize}

%----------------------------------------------------------------------------------------
% SYSTEM ARCHITECTURE
%----------------------------------------------------------------------------------------

\section{System Architecture}

\subsection{High-Level Architecture}

The Blockd platform employs a microservices architecture with three primary layers:

\begin{enumerate}
    \item \textbf{Interviewee Layer}: Custom Chromium browser with embedded security and eye tracking
    \item \textbf{Backend Layer}: 8 microservices handling authentication, sessions, AI analysis, and video
    \item \textbf{Interviewer Layer}: React web application for real-time monitoring and session management
\end{enumerate}

\subsection{Component Diagram}

\begin{figure}[H]
\centering
\begin{verbatim}
┌─────────────────────────────────────────────────────────────┐
│               INTERVIEWEE EXPERIENCE                        │
│  ┌───────────────────────────────────────────────────────┐  │
│  │      Blocked Chromium Browser (Desktop Native)        │  │
│  │  • Security Monitoring    • Eye Tracking              │  │
│  │  • Video Capture          • System Telemetry          │  │
│  │  • Fullscreen Lock        • Process Detection         │  │
│  └───────────────────────────────────────────────────────┘  │
└─────────────────────────────────────────────────────────────┘
                        │
                        │ WebSocket + Protocol Buffers
                        ▼
┌─────────────────────────────────────────────────────────────┐
│                  BACKEND MICROSERVICES                       │
│  ┌─────────────┐  ┌─────────────┐  ┌─────────────┐        │
│  │ API Gateway │  │   Video     │  │     AI      │        │
│  │  (Fastify)  │  │  Service    │  │  Detection  │        │
│  └─────────────┘  └─────────────┘  └─────────────┘        │
│  ┌─────────────┐  ┌─────────────┐  ┌─────────────┐        │
│  │     Eye     │  │  Response   │  │    Auth     │        │
│  │  Tracking   │  │   Timing    │  │  Service    │        │
│  └─────────────┘  └─────────────┘  └─────────────┘        │
└─────────────────────────────────────────────────────────────┘
                        │
                        │ HTTPS + WebSocket
                        ▼
┌─────────────────────────────────────────────────────────────┐
│                  INTERVIEWER EXPERIENCE                      │
│  ┌───────────────────────────────────────────────────────┐  │
│  │        React Web App (Browser-based)                  │  │
│  │  • Real-time Video       • Security Dashboard         │  │
│  │  • Session Management    • AI Detection Results       │  │
│  │  • Gaze Heatmaps         • Comprehensive Reports      │  │
│  └───────────────────────────────────────────────────────┘  │
└─────────────────────────────────────────────────────────────┘
\end{verbatim}
\caption{Blockd Platform Architecture}
\end{figure}

\subsection{Data Flow Architecture}

\begin{figure}[H]
\centering
\begin{verbatim}
Interview Question Asked
         │
         ├─► AI Detection Service
         │   ├─► Generate AI answers (GPT-4, Claude, Gemini)
         │   ├─► Compute embeddings (pgvector)
         │   └─► Cache in Redis (24h TTL)
         │
Interviewee Responds (Audio/Text)
         │
         ├─► Speech-to-Text (Whisper API)
         ├─► Response Timing Analysis
         ├─► AI Similarity Scoring
         │   ├─► Semantic similarity (cosine, 384-dim vectors)
         │   ├─► N-gram overlap (Jaccard index)
         │   ├─► Perplexity scoring (GPT-2)
         │   ├─► Stylometric analysis
         │   └─► XGBoost ensemble → Risk Score (0.0-1.0)
         │
Continuous Monitoring (During Session)
         │
         ├─► Eye Tracking (30 FPS)
         │   ├─► MediaPipe FaceMesh (468 facial landmarks)
         │   ├─► Gaze vector computation
         │   ├─► Kalman filtering (smooth tracking)
         │   ├─► Off-screen detection
         │   └─► Pattern recognition (reading, drift)
         │
         ├─► Security Telemetry (every 5 seconds)
         │   ├─► Process enumeration (detect OBS, VMware, etc.)
         │   ├─► Window focus tracking
         │   ├─► Screen recording detection
         │   ├─► VM detection (CPUID, registry, SMBIOS)
         │   ├─► Clipboard monitoring
         │   └─► System metrics (CPU, memory)
         │
         └─► Store in TimescaleDB
                     │ (30-day retention, 7-day compression)
                     ▼
         Generate Risk Score & Comprehensive Report (PDF)
\end{verbatim}
\caption{Blockd Data Flow}
\end{figure}

%----------------------------------------------------------------------------------------
% TECHNOLOGY STACK
%----------------------------------------------------------------------------------------

\section{Technology Stack}

\subsection{Frontend Technologies}

\subsubsection{Interviewer Web Application}

The interviewer-facing application is a modern single-page application built with cutting-edge web technologies.

\begin{table}[H]
\centering
\begin{tabularx}{\textwidth}{lX}
\toprule
\textbf{Technology} & \textbf{Purpose \& Version} \\
\midrule
React & v19.2.0 - UI framework with concurrent features \\
TypeScript & v5.9.3 - Type-safe development \\
Vite & v6.x - Fast build tool with HMR \\
Tailwind CSS & v4.x - Utility-first CSS framework \\
shadcn/ui & Latest - Pre-built accessible components \\
Zustand & v5.x - Lightweight state management \\
TanStack Query & v5.x - Server state management \\
React Router & v7.x - Client-side routing \\
Socket.io Client & v4.x - WebSocket real-time communication \\
mediasoup-client & v3.x - WebRTC video streaming \\
React Hook Form & Latest - Form state management \\
Zod & Latest - Runtime type validation \\
\bottomrule
\end{tabularx}
\caption{Frontend Technology Stack}
\end{table}

\subsubsection{Interviewee Chromium Browser}

A complete fork of Chromium with custom security and monitoring capabilities.

\begin{itemize}[noitemsep]
    \item \textbf{Base Version}: Chromium 142.0.7444.175 (Stable, November 2025)
    \item \textbf{Build System}: GN + Ninja build system
    \item \textbf{Languages}: C++ (browser core), JavaScript/TypeScript (UI), Objective-C++ (macOS)
    \item \textbf{Build Requirements}: 100+ GB disk, 16+ GB RAM, 8+ CPU cores
    \item \textbf{Build Time}: 2-8 hours initial, 10-60 minutes incremental
\end{itemize}

Key modifications to Chromium source:

\begin{lstlisting}[language=C++, caption=Chromium Modifications]
// Browser Process
chrome/browser/blocked/blocked_security/     // Security monitoring
chrome/browser/blocked/blocked_telemetry/    // System telemetry
chrome/browser/blocked/blocked_ipc/          // Backend communication
chrome/browser/blocked/blocked_video/        // Video capture
chrome/browser/ui/blocked/                   // UI restrictions

// Renderer Process
content/renderer/blocked_eye_tracking/       // Eye tracking module
content/renderer/blocked_video/              // Video frame processing
content/renderer/blocked_ipc/                // Renderer-to-browser IPC

// Third-party Integration
third_party/mediapipe/                       // MediaPipe FaceMesh
\end{lstlisting}

\subsection{Backend Technologies}

\subsubsection{Microservices Stack}

\begin{table}[H]
\centering
\small
\begin{tabularx}{\textwidth}{lXX}
\toprule
\textbf{Service} & \textbf{Technology} & \textbf{Purpose} \\
\midrule
API Gateway & Fastify 5.x, Node.js 24.11.0 & Request routing, rate limiting, JWT auth \\
Auth Service & Node.js 24.11.0, bcrypt & Authentication, MFA, OAuth 2.0 \\
Session Service & Node.js 24.11.0, Prisma & Session lifecycle, state machine \\
WebSocket Service & Socket.io 4.x, Redis adapter & Real-time bidirectional communication \\
AI Detection & Python 3.14.0, PyTorch 2.5.x & Multi-LLM comparison, XGBoost classification \\
Eye Tracking & Python 3.14.0, MediaPipe 0.10.x & Gaze analysis, LSTM anomaly detection \\
Response Timing & Python 3.14.0, Whisper & Speech-to-text, timing anomaly detection \\
Video Service & Python 3.14.0, mediasoup 3.19.11 & WebRTC SFU, FFmpeg recording \\
\bottomrule
\end{tabularx}
\caption{Microservices Technology Stack}
\end{table}

\subsubsection{Database Layer}

\begin{table}[H]
\centering
\begin{tabularx}{\textwidth}{lXl}
\toprule
\textbf{Database} & \textbf{Purpose} & \textbf{Version} \\
\midrule
PostgreSQL & Primary relational data (users, sessions) & 18.1 \\
TimescaleDB & Time-series data (gaze, telemetry, events) & 2.x \\
pgvector & Vector embeddings for AI similarity & 0.7.x \\
Redis Cluster & Caching, sessions, rate limiting & 8.4 \\
RabbitMQ & Message queue, async job processing & 4.x \\
\bottomrule
\end{tabularx}
\caption{Database Technologies}
\end{table}

Database schema highlights:

\begin{itemize}[noitemsep]
    \item \textbf{11 core tables}: users, organizations, sessions, questions, answers, etc.
    \item \textbf{31 indexes}: Optimized for query performance
    \item \textbf{2 hypertables}: gaze\_events, browser\_telemetry (TimescaleDB)
    \item \textbf{Retention policies}: 30-day default, 7-day compression
    \item \textbf{Replication}: 1 primary + 2 read replicas
\end{itemize}

\subsection{Infrastructure Technologies}

\begin{table}[H]
\centering
\begin{tabularx}{\textwidth}{lX}
\toprule
\textbf{Technology} & \textbf{Purpose} \\
\midrule
Docker & Containerization (8 multi-stage Dockerfiles) \\
Kubernetes & Orchestration (v1.31, 15+ manifests) \\
Helm & Package management (charts for dev/staging/prod) \\
Terraform & Infrastructure as Code (4 modules: cluster, database, storage, monitoring) \\
GitHub Actions & CI/CD pipelines (5 workflows: test, build, deploy-staging, deploy-prod, security) \\
Prometheus & Metrics collection (16 production alerts) \\
Grafana & Visualization (6 pre-built dashboards) \\
External Secrets Operator & Secrets management (AWS Secrets Manager integration) \\
\bottomrule
\end{tabularx}
\caption{Infrastructure Technologies}
\end{table}

%----------------------------------------------------------------------------------------
% COMPONENT SPECIFICATIONS
%----------------------------------------------------------------------------------------

\section{Component Specifications}

\subsection{Database Layer}

\subsubsection{PostgreSQL Schema}

The database schema consists of 11 core tables supporting the complete interview workflow:

\begin{lstlisting}[language=SQL, caption=Core Tables]
-- Users and Organizations
CREATE TABLE organizations (
  id UUID PRIMARY KEY DEFAULT uuid_generate_v4(),
  name VARCHAR(255) NOT NULL,
  created_at TIMESTAMP DEFAULT NOW()
);

CREATE TABLE users (
  id UUID PRIMARY KEY DEFAULT uuid_generate_v4(),
  email VARCHAR(255) UNIQUE NOT NULL,
  password_hash VARCHAR(255),  -- bcrypt (12 rounds)
  full_name VARCHAR(255) NOT NULL,
  role user_role NOT NULL,     -- admin, interviewer, interviewee, viewer
  organization_id UUID REFERENCES organizations(id),
  mfa_enabled BOOLEAN DEFAULT FALSE,
  mfa_secret VARCHAR(32),      -- Base32-encoded TOTP secret
  email_verified BOOLEAN DEFAULT FALSE,
  created_at TIMESTAMP DEFAULT NOW(),
  updated_at TIMESTAMP DEFAULT NOW()
);

-- Sessions and Interviews
CREATE TABLE interview_sessions (
  id UUID PRIMARY KEY DEFAULT uuid_generate_v4(),
  organization_id UUID REFERENCES organizations(id),
  interviewer_id UUID REFERENCES users(id),
  interviewee_id UUID REFERENCES users(id),
  status session_status NOT NULL,
  scheduled_start TIMESTAMP,
  actual_start TIMESTAMP,
  actual_end TIMESTAMP,
  session_token VARCHAR(255) UNIQUE,
  video_recording_url TEXT,
  risk_score DECIMAL(3,2),     -- 0.00 to 1.00
  created_at TIMESTAMP DEFAULT NOW(),
  updated_at TIMESTAMP DEFAULT NOW()
);

-- Security Events
CREATE TABLE security_events (
  id UUID PRIMARY KEY DEFAULT uuid_generate_v4(),
  session_id UUID REFERENCES interview_sessions(id),
  event_type security_event_type NOT NULL,
  severity event_severity NOT NULL,
  description TEXT,
  metadata JSONB,
  timestamp TIMESTAMP DEFAULT NOW()
);

-- AI Answer Cache
CREATE TABLE ai_answer_cache (
  id UUID PRIMARY KEY DEFAULT uuid_generate_v4(),
  question_hash VARCHAR(64) NOT NULL,
  model_name VARCHAR(50) NOT NULL,
  answer_text TEXT NOT NULL,
  embedding vector(384),       -- pgvector
  perplexity_score FLOAT,
  generated_at TIMESTAMP DEFAULT NOW(),
  access_count INTEGER DEFAULT 0,
  last_accessed TIMESTAMP,
  UNIQUE(question_hash, model_name)
);

-- TimescaleDB Hypertables
CREATE TABLE gaze_events (
  session_id UUID NOT NULL,
  timestamp TIMESTAMP NOT NULL,
  gaze_x FLOAT,                -- Normalized 0-1
  gaze_y FLOAT,                -- Normalized 0-1
  is_off_screen BOOLEAN,
  off_screen_direction VARCHAR(20),
  confidence FLOAT
);

SELECT create_hypertable('gaze_events', 'timestamp');
SELECT add_retention_policy('gaze_events', INTERVAL '30 days');
SELECT add_compression_policy('gaze_events', INTERVAL '7 days');
\end{lstlisting}

\subsubsection{Redis Configuration}

Redis cluster configuration for high availability:

\begin{itemize}[noitemsep]
    \item \textbf{Topology}: 6 nodes (3 primaries + 3 replicas)
    \item \textbf{Memory}: 64 GB total (21.3 GB per primary)
    \item \textbf{Persistence}: AOF with fsync every second
    \item \textbf{Eviction}: allkeys-lru for cache management
    \item \textbf{Use Cases}:
    \begin{itemize}
        \item Session tokens (1h TTL)
        \item Rate limiting counters
        \item AI answer cache (24h TTL)
        \item Real-time session data
        \item WebSocket connection state
    \end{itemize}
\end{itemize}

\subsubsection{RabbitMQ Topology}

Message queue architecture for asynchronous processing:

\begin{lstlisting}[caption=RabbitMQ Exchanges and Queues]
Exchanges:
1. video_processing (topic) - Video encoding and upload
2. ai_detection (direct) - AI analysis requests
3. security_events (fanout) - Real-time security alerts
4. gaze_analysis (direct) - Eye tracking pattern analysis
5. dlx (dead letter exchange) - Failed message handling

Queues (20 total):
- video_processing.encode, video_processing.upload
- ai_detection.similarity, ai_detection.perplexity
- gaze_analysis.pattern, gaze_analysis.anomaly
- security_events.high_priority, security_events.normal
- [12 more specialized queues...]

Workers:
- Celery workers (Python) for task execution
- 10-50 workers per queue (auto-scaling)
- Max retries: 3, Retry delay: exponential backoff
\end{lstlisting}

\subsection{Backend Microservices}

\subsubsection{API Gateway}

The API Gateway serves as the single entry point for all client requests, handling authentication, rate limiting, and request routing.

\begin{lstlisting}[language=JavaScript, caption=API Gateway Architecture]
// Fastify 5.x with TypeScript
import Fastify from 'fastify';
import { fastifyJWT } from '@fastify/jwt';

const app = Fastify({ logger: true });

// JWT Authentication
app.register(fastifyJWT, {
  secret: { public: publicKey, private: privateKey },
  sign: { algorithm: 'RS256', expiresIn: '1h' }
});

// Rate Limiting (Redis-backed)
app.register(rateLimit, {
  max: 100,              // 100 requests
  timeWindow: '1 minute', // per minute
  redis: redisClient
});

// Routes (30+ endpoints)
app.register(authRoutes, { prefix: '/api/v1/auth' });
app.register(sessionRoutes, { prefix: '/api/v1/sessions' });
app.register(analysisRoutes, { prefix: '/api/v1/analysis' });
app.register(browserRoutes, { prefix: '/api/v1/browser' });

// Error Handling
app.setErrorHandler((error, request, reply) => {
  logger.error({ error, request });
  reply.status(500).send({ error: 'Internal Server Error' });
});
\end{lstlisting}

Key endpoints:

\begin{table}[H]
\centering
\small
\begin{tabularx}{\textwidth}{llX}
\toprule
\textbf{Method} & \textbf{Endpoint} & \textbf{Description} \\
\midrule
POST & /api/v1/auth/register & User registration \\
POST & /api/v1/auth/login & Authentication with JWT \\
POST & /api/v1/auth/refresh & Refresh access token \\
POST & /api/v1/sessions & Create interview session \\
GET & /api/v1/sessions/:id & Get session details \\
POST & /api/v1/sessions/:id/start & Start session \\
POST & /api/v1/sessions/:id/end & End session \\
POST & /api/v1/analysis/answer & Analyze answer with AI \\
POST & /api/v1/browser/security/event & Report security event \\
POST & /api/v1/browser/telemetry/batch & Batch upload telemetry \\
\bottomrule
\end{tabularx}
\caption{Key API Endpoints}
\end{table}

\subsubsection{Authentication Service}

Handles user authentication with multiple methods and MFA support.

Features:
\begin{itemize}[noitemsep]
    \item \textbf{JWT Tokens}: RS256 algorithm, 1h access token, 7d refresh token
    \item \textbf{MFA}: TOTP (Time-based One-Time Password) with QR code generation
    \item \textbf{OAuth 2.0}: Google and Microsoft SSO integration
    \item \textbf{Password Security}: bcrypt with 12 rounds, minimum 12 characters
    \item \textbf{Session Management}: Redis-backed with automatic expiration
    \item \textbf{Role-Based Access Control}: admin, interviewer, interviewee, viewer
\end{itemize}

\subsubsection{AI Detection Service}

The AI Detection Service is the core intellectual property of the platform, implementing sophisticated analysis to detect AI-generated answers.

\begin{lstlisting}[language=Python, caption=AI Detection Pipeline]
class AIDetectionService:
    """Multi-LLM comparison with XGBoost ensemble classification"""

    def __init__(self):
        self.llm_clients = {
            'gpt-4': OpenAIClient(),
            'claude-3.5-sonnet': AnthropicClient(),
            'gemini-1.5-pro': GoogleClient()
        }
        self.embedding_model = SentenceTransformer('all-MiniLM-L6-v2')
        self.xgboost_model = XGBClassifier.load('ai_detector_v2.model')

    async def analyze_answer(self, question: str, answer: str) -> AnalysisResult:
        # 1. Generate AI answers from all LLMs (with caching)
        ai_answers = await self.generate_ai_answers(question)

        # 2. Compute embeddings (384-dimensional vectors)
        human_embedding = self.embedding_model.encode(answer)
        ai_embeddings = {
            model: self.embedding_model.encode(ai_ans)
            for model, ai_ans in ai_answers.items()
        }

        # 3. Calculate semantic similarity (cosine)
        similarities = {
            model: cosine_similarity(human_embedding, ai_emb)
            for model, ai_emb in ai_embeddings.items()
        }

        # 4. Calculate perplexity score using GPT-2
        perplexity = self.calculate_perplexity(answer)

        # 5. N-gram overlap (Jaccard index)
        ngram_scores = self.calculate_ngram_overlap(answer, ai_answers)

        # 6. Stylometric analysis
        style_features = self.extract_style_features(answer)

        # 7. XGBoost ensemble prediction
        features = self.combine_features(
            similarities, perplexity, ngram_scores, style_features
        )
        risk_score = self.xgboost_model.predict_proba(features)[0][1]

        # 8. Generate flags
        flags = self.generate_flags(similarities, perplexity, risk_score)

        return AnalysisResult(
            risk_score=risk_score,
            risk_level=self.categorize_risk(risk_score),
            similarities=similarities,
            perplexity_score=perplexity,
            flags=flags
        )
\end{lstlisting}

AI Detection Risk Categorization:

\begin{table}[H]
\centering
\begin{tabularx}{\textwidth}{llX}
\toprule
\textbf{Risk Level} & \textbf{Score Range} & \textbf{Interpretation} \\
\midrule
Minimal & 0.00 - 0.49 & Likely human-generated \\
Low & 0.50 - 0.69 & Probably human, minor similarities \\
Medium & 0.70 - 0.84 & Investigate further, notable AI patterns \\
High & 0.85 - 1.00 & Likely AI-generated, strong evidence \\
\bottomrule
\end{tabularx}
\caption{AI Detection Risk Levels}
\end{table}

\subsubsection{Eye Tracking Service}

Real-time gaze analysis using computer vision and machine learning.

Architecture:
\begin{itemize}[noitemsep]
    \item \textbf{Face Detection}: MediaPipe FaceMesh with 468 facial landmarks
    \item \textbf{Gaze Estimation}: Custom CNN based on ResNet-18 architecture
    \item \textbf{Kalman Filtering}: Smooth tracking with noise reduction
    \item \textbf{LSTM Autoencoder}: Anomaly detection for suspicious patterns
    \item \textbf{Performance}: 30 FPS, <50ms latency, <15\% CPU usage
\end{itemize}

Gaze analysis pipeline:

\begin{lstlisting}[language=Python, caption=Eye Tracking Pipeline]
class GazeAnalysisService:
    """30 FPS real-time gaze tracking and analysis"""

    def process_frame(self, frame: np.ndarray) -> GazeData:
        # 1. Detect face with MediaPipe (468 landmarks)
        landmarks = self.facemesh.process(frame)

        # 2. Extract eye/iris landmarks
        left_eye = landmarks[LEFT_EYE_INDICES]
        right_eye = landmarks[RIGHT_EYE_INDICES]

        # 3. Estimate gaze vector (3D)
        gaze_vector = self.gaze_estimator.predict(left_eye, right_eye)

        # 4. Apply Kalman filter for smoothing
        filtered_gaze = self.kalman_filter.update(gaze_vector)

        # 5. Apply calibration matrix (3x3 transform)
        calibrated_gaze = self.calibration_matrix @ filtered_gaze

        # 6. Convert to screen coordinates (0-1 normalized)
        screen_x, screen_y = self.to_screen_coords(calibrated_gaze)

        # 7. Detect off-screen gaze
        is_off_screen, direction = self.detect_off_screen(
            screen_x, screen_y
        )

        # 8. Pattern analysis with LSTM
        pattern_score = self.lstm_anomaly_detector.score(gaze_sequence)

        return GazeData(
            x=screen_x,
            y=screen_y,
            is_off_screen=is_off_screen,
            direction=direction,
            confidence=landmarks.confidence,
            anomaly_score=pattern_score
        )
\end{lstlisting}

\subsubsection{Video Service}

WebRTC-based video streaming with mediasoup SFU (Selective Forwarding Unit).

\begin{itemize}[noitemsep]
    \item \textbf{Technology}: mediasoup 3.19.11 (WebRTC SFU)
    \item \textbf{Capacity}: 100+ concurrent streams
    \item \textbf{Resolution}: Adaptive (240p - 720p)
    \item \textbf{Frame Rate}: 30 FPS default
    \item \textbf{Recording}: FFmpeg with hardware acceleration (NVENC, QuickSync)
    \item \textbf{Storage}: S3-compatible (AWS S3, MinIO, Cloudflare R2)
    \item \textbf{Latency}: <1 second end-to-end
\end{itemize}

\subsection{Frontend Application}

\subsubsection{Architecture Overview}

The interviewer web application is a modern React SPA with the following structure:

\begin{lstlisting}[caption=Frontend Directory Structure]
frontend/interviewer-app/
├── src/
│   ├── components/          # UI components (shadcn/ui + custom)
│   │   ├── ui/             # 20+ shadcn/ui primitives
│   │   ├── VideoPlayer.tsx # WebRTC video display
│   │   ├── SecurityEventsDashboard.tsx
│   │   ├── GazeHeatmap.tsx
│   │   ├── AIDetectionResults.tsx
│   │   └── RealtimeChat.tsx
│   ├── pages/              # Route pages
│   │   ├── LoginPage.tsx
│   │   ├── DashboardPage.tsx
│   │   ├── SessionsPage.tsx
│   │   ├── SessionDetailPage.tsx
│   │   └── SettingsPage.tsx
│   ├── hooks/              # Custom React hooks
│   │   ├── useAuth.ts
│   │   ├── useWebSocket.ts
│   │   ├── useWebRTC.ts
│   │   └── useSession.ts
│   ├── stores/             # Zustand state management
│   │   ├── auth-store.ts
│   │   ├── session-store.ts
│   │   └── realtime-store.ts
│   └── lib/                # Utilities
│       ├── api-client.ts
│       ├── validations.ts
│       └── utils.ts
├── package.json
├── vite.config.ts
└── tsconfig.json
\end{lstlisting}

\subsubsection{Key Features}

\textbf{Authentication \& Session Management:}
\begin{itemize}[noitemsep]
    \item Login with email/password or OAuth (Google, Microsoft)
    \item MFA support with TOTP
    \item Session creation with dynamic questions editor
    \item Session list with filters, search, and pagination
    \item Settings page with profile, security, and preferences
\end{itemize}

\textbf{Real-time Features:}
\begin{itemize}[noitemsep]
    \item WebSocket connection for real-time events
    \item WebRTC video player with controls
    \item Security events dashboard with live updates
    \item Gaze heatmap with canvas rendering
    \item AI detection results visualization
    \item Real-time chat with typing indicators
\end{itemize}

\subsection{Chromium Browser}

\subsubsection{Browser Process Security}

The browser process implements comprehensive security monitoring that runs with elevated privileges.

\begin{lstlisting}[language=C++, caption=Security Monitoring Service]
// chrome/browser/blocked/blocked_security/blocked_security_service.h
class BlockedSecurityService : public KeyedService {
 public:
  explicit BlockedSecurityService();
  ~BlockedSecurityService() override;

  void StartMonitoring();
  void StopMonitoring();

  // Observer pattern for event notification
  void AddObserver(SecurityEventObserver* observer);
  void RemoveObserver(SecurityEventObserver* observer);

 private:
  void OnMonitoringTick();
  void CheckProcesses();
  void CheckVirtualMachine();
  void CheckWindowFocus();
  void CheckClipboard();

  // Platform-specific monitors (Windows/macOS/Linux)
  std::unique_ptr<PlatformSecurityMonitor> platform_monitor_;

  // Risk scoring
  float CalculateRiskScore() const;

  // Event buffer (circular, 100 events max)
  std::deque<SecurityEvent> recent_events_;

  // Observers
  base::ObserverList<SecurityEventObserver> observers_;

  // Monitoring interval (5 seconds)
  base::RepeatingTimer monitoring_timer_;

  base::WeakPtrFactory<BlockedSecurityService> weak_factory_{this};
};
\end{lstlisting}

Platform-specific detection capabilities:

\begin{table}[H]
\centering
\small
\begin{tabularx}{\textwidth}{lXXX}
\toprule
\textbf{Detection} & \textbf{Windows} & \textbf{macOS} & \textbf{Linux} \\
\midrule
Screen Recorders & OBS, Camtasia, Bandicam & QuickTime, ScreenFlow, OBS & OBS, SimpleScreenRecorder, Kazam \\
VMs & VMware, VirtualBox, Hyper-V & Parallels, VMware Fusion, VirtualBox & QEMU/KVM, VirtualBox, VMware \\
Remote Access & TeamViewer, AnyDesk & TeamViewer, AnyDesk & TeamViewer, AnyDesk \\
Method & Win32 API, CPUID, Registry & IOKit, sysctl, NSRunningApp & /proc, DMI, kernel modules \\
\bottomrule
\end{tabularx}
\caption{Platform-Specific Security Detection}
\end{table}

\subsubsection{Renderer Process Eye Tracking}

Eye tracking runs in the sandboxed renderer process, sending data to the browser process via Mojo IPC.

\begin{lstlisting}[language=C++, caption=Eye Tracking Pipeline]
// content/renderer/blocked_eye_tracking/eye_tracker.h
class EyeTracker {
 public:
  EyeTracker();
  ~EyeTracker();

  // Start tracking at 30 FPS
  void Start();
  void Stop();

  // Run calibration (9-point grid, 18 seconds)
  void Calibrate(CalibrationCallback callback);

 private:
  // Processing pipeline (runs on background thread)
  void ProcessFrame(const VideoFrame& frame);

  // MediaPipe FaceMesh integration
  std::unique_ptr<FaceDetector> face_detector_;

  // Gaze estimation with Kalman filtering
  std::unique_ptr<GazeEstimator> gaze_estimator_;

  // Background processing thread
  std::unique_ptr<base::Thread> worker_thread_;

  // Mojo IPC to browser process
  mojo::Remote<mojom::EyeTrackingHost> host_;

  // Calibration state (3x3 transformation matrix)
  Eigen::Matrix3f calibration_matrix_;

  // Performance: 30 FPS target
  static constexpr int kTargetFps = 30;
  static constexpr int kFrameIntervalMs = 1000 / kTargetFps;
};
\end{lstlisting}

Calibration process:

\begin{enumerate}
    \item Display 9 calibration points in 3x3 grid
    \item User looks at each point for 2 seconds
    \item Collect 60 gaze samples per point (30 FPS × 2 seconds)
    \item Calculate 3x3 transformation matrix using least squares
    \item Validate accuracy (target: >90\%)
    \item Store calibration in local state
    \item Apply to all future gaze estimates
\end{enumerate}

\subsubsection{Platform Distribution}

Complete installer and auto-update infrastructure for all platforms.

\begin{table}[H]
\centering
\begin{tabularx}{\textwidth}{lXXl}
\toprule
\textbf{Platform} & \textbf{Installer} & \textbf{Auto-Update} & \textbf{Signing} \\
\midrule
Windows & NSIS (.exe) & Google Omaha & Authenticode \\
macOS & DMG + App Bundle & Sparkle Framework & Developer ID + Notarization \\
Linux & DEB, RPM, AppImage & Custom HTTP & GPG \\
\bottomrule
\end{tabularx}
\caption{Platform Distribution}
\end{table}

%----------------------------------------------------------------------------------------
% TESTING & CI/CD
%----------------------------------------------------------------------------------------

\section{Testing \& CI/CD}

\subsection{Testing Infrastructure}

\subsubsection{Test Coverage}

\begin{table}[H]
\centering
\begin{tabularx}{\textwidth}{lXll}
\toprule
\textbf{Test Type} & \textbf{Tool} & \textbf{Files} & \textbf{Tests} \\
\midrule
E2E Tests & Playwright & 6 spec files & 43 scenarios \\
API Tests & Newman/Postman & 4 collections & 22 tests \\
Load Tests & k6 & 3 scenarios & 3 scripts \\
\textbf{Total} & & \textbf{13 files} & \textbf{68+ tests} \\
\bottomrule
\end{tabularx}
\caption{Test Coverage Summary}
\end{table}

\subsubsection{E2E Test Scenarios}

Playwright tests covering critical user journeys:

\begin{lstlisting}[caption=E2E Test Categories]
1. interview-flow.spec.ts
   - Complete interview lifecycle (create, join, conduct, end)

2. auth.spec.ts
   - Login with email/password
   - Login with MFA
   - OAuth flows (Google, Microsoft)
   - Registration with email verification

3. realtime.spec.ts
   - WebSocket connection and reconnection
   - Real-time security event updates
   - Gaze data streaming
   - Chat messaging

4. security.spec.ts
   - Process detection simulation
   - VM detection scenarios
   - Window focus tracking
   - Alert generation

5. ai-detection.spec.ts
   - Submit human answer
   - Submit AI-generated answer
   - Risk score validation
   - Flag generation

6. video.spec.ts
   - WebRTC connection establishment
   - Video quality adaptation
   - Recording start/stop
   - Multiple concurrent streams
\end{lstlisting}

\subsubsection{Load Testing}

k6 scenarios for performance validation:

\begin{table}[H]
\centering
\begin{tabularx}{\textwidth}{lXll}
\toprule
\textbf{Scenario} & \textbf{Configuration} & \textbf{Target} & \textbf{Duration} \\
\midrule
Normal Load & 100 VUs, steady state & 500 req/min & 10 minutes \\
Peak Load & 500 VUs, steady state & 2000 req/min & 5 minutes \\
Stress Test & 0→1000 VUs ramp & Find breaking point & 10 minutes \\
\bottomrule
\end{tabularx}
\caption{Load Testing Scenarios}
\end{table}

\subsection{CI/CD Automation}

\subsubsection{GitHub Actions Workflows}

\begin{table}[H]
\centering
\small
\begin{tabularx}{\textwidth}{lX}
\toprule
\textbf{Workflow} & \textbf{Stages} \\
\midrule
test.yml & Lint → Unit Tests (matrix: 9 services) → Integration Tests → E2E Tests \\
build.yml & Docker Build (matrix: 9 services) → Trivy Scan → Push to Registry \\
deploy-staging.yml & Deploy with Helm → Smoke Tests → Newman API Tests \\
deploy-production.yml & Create Green → Smoke Tests → Switch Traffic → Monitor 10min → Cleanup Blue \\
security.yml & Dependency Scan → Code Scan (CodeQL, SonarQube) → Container Scan \\
\bottomrule
\end{tabularx}
\caption{CI/CD Workflows}
\end{table}

\subsubsection{Deployment Strategy}

\textbf{Staging Deployment}:
\begin{itemize}[noitemsep]
    \item Trigger: Automatic on push to \texttt{develop} branch
    \item Strategy: Rolling update (max unavailable: 25\%)
    \item Testing: Automated smoke tests + API integration tests
    \item Rollback: Automatic on test failure
\end{itemize}

\textbf{Production Deployment}:
\begin{itemize}[noitemsep]
    \item Trigger: Manual approval + push to \texttt{main} branch
    \item Strategy: Blue-green deployment (zero downtime)
    \item Process:
    \begin{enumerate}
        \item Deploy green environment (parallel to blue)
        \item Run comprehensive smoke tests
        \item Switch traffic from blue to green
        \item Monitor for 10 minutes (CPU, memory, errors, latency)
        \item Automatic rollback if health checks fail
        \item Cleanup blue environment after 24 hours
    \end{enumerate}
\end{itemize}

\subsubsection{Monitoring \& Alerting}

\textbf{Prometheus Alerts (16 total)}:

\begin{table}[H]
\centering
\small
\begin{tabularx}{\textwidth}{lXl}
\toprule
\textbf{Alert} & \textbf{Condition} & \textbf{Severity} \\
\midrule
HighErrorRate & Error rate > 1\% for 5 minutes & Critical \\
HighLatency & p95 latency > 500ms for 5 minutes & Warning \\
ServiceDown & Service unavailable for 1 minute & Critical \\
DatabasePoolExhausted & Available connections < 10 & Critical \\
RedisMemoryHigh & Memory usage > 80\% & Warning \\
PodCPUHigh & CPU usage > 80\% for 10 minutes & Warning \\
DiskSpaceLow & Disk usage > 85\% & Warning \\
\bottomrule
\end{tabularx}
\caption{Key Prometheus Alerts}
\end{table}

%----------------------------------------------------------------------------------------
% DEPLOYMENT GUIDE
%----------------------------------------------------------------------------------------

\section{Deployment Guide}

\subsection{Prerequisites}

\begin{enumerate}
    \item \textbf{Kubernetes Cluster}
    \begin{itemize}
        \item AWS EKS, Google GKE, or Azure AKS recommended
        \item Minimum: 3 nodes (c6i.2xlarge or equivalent)
        \item Kubernetes version: 1.31+
    \end{itemize}

    \item \textbf{Domain \& DNS}
    \begin{itemize}
        \item Domain name (e.g., blockd.com)
        \item DNS management (Route 53, CloudFlare, etc.)
        \item SSL/TLS certificates (Let's Encrypt via cert-manager)
    \end{itemize}

    \item \textbf{External Services}
    \begin{itemize}
        \item AWS S3 or compatible object storage
        \item LLM API keys (OpenAI, Anthropic, Google)
        \item SMTP server for email notifications
        \item OAuth credentials (Google, Microsoft)
    \end{itemize}

    \item \textbf{Code Signing Certificates}
    \begin{itemize}
        \item Windows: EV Code Signing Certificate
        \item macOS: Apple Developer ID Certificate
        \item Linux: GPG key for package signing
    \end{itemize}
\end{enumerate}

\subsection{Infrastructure Provisioning}

\subsubsection{Terraform Setup}

\begin{lstlisting}[language=bash, caption=Provision Infrastructure with Terraform]
cd infrastructure/terraform

# Initialize Terraform
terraform init

# Create production infrastructure
terraform plan -var-file=environments/production/terraform.tfvars
terraform apply -var-file=environments/production/terraform.tfvars

# Outputs:
# - EKS cluster endpoint
# - RDS PostgreSQL endpoint
# - ElastiCache Redis endpoint
# - S3 bucket names
\end{lstlisting}

Terraform modules provision:
\begin{itemize}[noitemsep]
    \item EKS cluster with auto-scaling node groups
    \item RDS PostgreSQL 18 with Multi-AZ replication
    \item ElastiCache Redis 8.4 cluster (6 nodes)
    \item Amazon MQ RabbitMQ 4.x cluster
    \item S3 buckets with lifecycle policies
    \item CloudFront CDN distribution
    \item VPC with public/private subnets
    \item KMS keys for encryption
\end{itemize}

\subsection{Database Setup}

\begin{lstlisting}[language=bash, caption=Database Migration]
# Connect to PostgreSQL
export PGHOST=<rds-endpoint>
export PGDATABASE=blockd
export PGUSER=postgres
export PGPASSWORD=<password>

# Run Alembic migrations
cd database
alembic upgrade head

# Verify tables
psql -c "\dt"

# Create TimescaleDB hypertables
psql -f create_hypertables.sql

# Seed initial data (optional)
psql -f seed_data.sql
\end{lstlisting}

\subsection{Kubernetes Deployment}

\subsubsection{Install Dependencies}

\begin{lstlisting}[language=bash, caption=Install Kubernetes Dependencies]
# Install cert-manager for TLS certificates
kubectl apply -f https://github.com/cert-manager/cert-manager/releases/download/v1.13.0/cert-manager.yaml

# Install External Secrets Operator
kubectl apply -f k8s/external-secrets/external-secrets-operator.yaml

# Install Prometheus + Grafana
helm repo add prometheus-community https://prometheus-community.github.io/helm-charts
helm install monitoring prometheus-community/kube-prometheus-stack \
  -n monitoring --create-namespace

# Install Ingress NGINX
helm repo add ingress-nginx https://kubernetes.github.io/ingress-nginx
helm install ingress-nginx ingress-nginx/ingress-nginx \
  -n ingress-nginx --create-namespace
\end{lstlisting}

\subsubsection{Deploy Blockd Platform}

\begin{lstlisting}[language=bash, caption=Deploy with Helm]
# Create namespace
kubectl create namespace production

# Create secrets in AWS Secrets Manager
# (Database credentials, API keys, JWT keys, etc.)

# Configure External Secrets
kubectl apply -f k8s/external-secrets/ -n production

# Deploy Blockd platform
helm install blockd k8s/helm/blockd \
  -n production \
  -f k8s/helm/blockd/values-production.yaml

# Verify deployment
kubectl get pods -n production
kubectl get ingress -n production

# Check logs
kubectl logs -f deployment/api-gateway -n production
\end{lstlisting}

Helm chart deploys:
\begin{itemize}[noitemsep]
    \item 9 microservice deployments with HPA
    \item 9 Kubernetes services
    \item Ingress with TLS termination
    \item ConfigMaps and Secrets
    \item PersistentVolumeClaims for PostgreSQL
    \item NetworkPolicies for security
\end{itemize}

\subsection{Chromium Browser Build}

\begin{lstlisting}[language=bash, caption=Build Blocked Browser]
cd chromium

# Setup build environment (100+ GB disk, 16+ GB RAM)
./setup.sh

# This will:
# - Install depot_tools
# - Download Chromium source (~30 GB)
# - Install dependencies
# - Configure environment

# Apply Blocked modifications
./patches/apply-patches.sh

# Build browser (2-8 hours depending on CPU)
./build.sh --release

# Output:
# - Windows: out/Blocked/chrome.exe
# - macOS: out/Blocked/Chromium.app
# - Linux: out/Blocked/chrome

# Create installers
cd installer/windows
./build_installer.bat  # Windows NSIS installer

cd ../mac
./create_dmg.sh        # macOS DMG

cd ../linux
./build_appimage.sh    # Linux AppImage
\end{lstlisting}

\subsection{Post-Deployment Verification}

\begin{lstlisting}[language=bash, caption=Verification Checklist]
# Check all pods are running
kubectl get pods -n production

# Check ingress has external IP
kubectl get ingress -n production

# Test API Gateway
curl https://api.blockd.com/api/v1/health

# Test authentication
curl -X POST https://api.blockd.com/api/v1/auth/login \
  -H "Content-Type: application/json" \
  -d '{"email":"test@example.com","password":"password123"}'

# Run smoke tests
cd tests
npm run test:api -- --env production

# Check monitoring
kubectl port-forward -n monitoring svc/monitoring-grafana 3000:80
# Open http://localhost:3000

# Check logs for errors
kubectl logs -f deployment/api-gateway -n production | grep ERROR
\end{lstlisting}

%----------------------------------------------------------------------------------------
% SECURITY ARCHITECTURE
%----------------------------------------------------------------------------------------

\section{Security Architecture}

\subsection{Threat Model}

\subsubsection{Threats Mitigated}

\begin{table}[H]
\centering
\small
\begin{tabularx}{\textwidth}{lXX}
\toprule
\textbf{Threat} & \textbf{Detection Method} & \textbf{Mitigation} \\
\midrule
AI Assistance & Multi-LLM similarity, perplexity, n-gram & Real-time alert, risk score \\
Second Monitor & Eye tracking, off-screen detection & Log events, pattern analysis \\
Virtual Machine & CPUID, registry, SMBIOS checks & Block session start \\
Screen Recording & Process enumeration, window detection & Terminate process or session \\
Window Switching & Window focus monitoring & Log all focus changes \\
Clipboard Use & Clipboard monitoring & Log paste events \\
Remote Access & Process detection (TeamViewer, etc.) & Block session \\
\bottomrule
\end{tabularx}
\caption{Threat Detection and Mitigation}
\end{table}

\subsection{Authentication Security}

\begin{itemize}[noitemsep]
    \item \textbf{Password Hashing}: bcrypt with 12 rounds (cost factor)
    \item \textbf{JWT Tokens}: RS256 algorithm, 1-hour expiration
    \item \textbf{Refresh Tokens}: 7-day expiration with rotation
    \item \textbf{MFA}: TOTP with 30-second time step, 6-digit codes
    \item \textbf{OAuth 2.0}: Secure delegation with state parameter
    \item \textbf{Rate Limiting}: 100 requests/minute per IP (public), 500/minute (authenticated)
    \item \textbf{Session Security}: HttpOnly cookies, SameSite=Strict
\end{itemize}

\subsection{Data Security}

\begin{itemize}[noitemsep]
    \item \textbf{Encryption at Rest}: AES-256 for databases and S3
    \item \textbf{Encryption in Transit}: TLS 1.3 for all HTTPS traffic
    \item \textbf{PII Protection}: Database-level encryption, access logging
    \item \textbf{Video Retention}: Automatic deletion after 90 days
    \item \textbf{GDPR Compliance}: Data export and deletion workflows
    \item \textbf{Audit Logging}: All authentication events, session lifecycle, security events
\end{itemize}

\subsection{Network Security}

\begin{itemize}[noitemsep]
    \item \textbf{Network Policies}: Kubernetes NetworkPolicies restrict pod-to-pod traffic
    \item \textbf{Ingress Security}: TLS termination, rate limiting, WAF
    \item \textbf{Service Mesh}: Optional Istio for mTLS between services
    \item \textbf{Secrets Management}: External Secrets Operator with AWS Secrets Manager
    \item \textbf{IRSA}: IAM Roles for Service Accounts (pod-level IAM)
\end{itemize}

%----------------------------------------------------------------------------------------
% MONITORING & OPERATIONS
%----------------------------------------------------------------------------------------

\section{Monitoring \& Operations}

\subsection{Metrics Collection}

\textbf{Prometheus Configuration}:
\begin{itemize}[noitemsep]
    \item Scrape interval: 15 seconds
    \item Retention: 15 days
    \item Storage: 50 GB persistent volume
    \item Targets: All pods with \texttt{prometheus.io/scrape: "true"} annotation
\end{itemize}

\textbf{Key Metrics}:
\begin{table}[H]
\centering
\small
\begin{tabularx}{\textwidth}{lX}
\toprule
\textbf{Category} & \textbf{Metrics} \\
\midrule
HTTP & Request rate, latency (p50, p95, p99), error rate, status codes \\
Database & Connection pool usage, query time, active connections \\
Redis & Hit rate, memory usage, eviction rate, command latency \\
WebSocket & Active connections, messages/second, connection errors \\
Video & Active streams, bitrate, packet loss, encoding time \\
AI Detection & Analysis time, cache hit rate, LLM API latency \\
Eye Tracking & Frame rate, processing latency, off-screen events \\
\bottomrule
\end{tabularx}
\caption{Monitored Metrics}
\end{table}

\subsection{Grafana Dashboards}

Six pre-built dashboards for comprehensive visualization:

\begin{enumerate}
    \item \textbf{API Gateway}: Request rate, latency, error rate by endpoint
    \item \textbf{Database}: Connection pool, query time, replication lag
    \item \textbf{Redis}: Memory usage, hit rate, command distribution
    \item \textbf{WebSocket}: Connections over time, message throughput, errors
    \item \textbf{AI Detection}: Analysis time, cache hit rate, risk score distribution
    \item \textbf{Video Streaming}: Active streams, bitrate, quality metrics
\end{enumerate}

\subsection{Incident Response}

\subsubsection{Severity Levels}

\begin{table}[H]
\centering
\begin{tabularx}{\textwidth}{llXl}
\toprule
\textbf{SEV} & \textbf{Description} & \textbf{Examples} & \textbf{Response} \\
\midrule
SEV1 & Critical outage & Platform down, data breach & 15 min \\
SEV2 & Major impact & Service degraded, 20\% error rate & 1 hour \\
SEV3 & Minor impact & Non-critical feature broken & 4 hours \\
SEV4 & Low impact & Visual bug, minor inconvenience & 24 hours \\
\bottomrule
\end{tabularx}
\caption{Incident Severity Levels}
\end{table}

\subsubsection{Rollback Procedure}

\begin{lstlisting}[language=bash, caption=Emergency Rollback]
# Rollback to previous Helm release
helm rollback blockd -n production

# Or rollback specific deployment
kubectl rollout undo deployment/api-gateway -n production

# Verify rollback
kubectl rollout status deployment/api-gateway -n production

# Check logs
kubectl logs -f deployment/api-gateway -n production
\end{lstlisting}

%----------------------------------------------------------------------------------------
% PERFORMANCE BENCHMARKS
%----------------------------------------------------------------------------------------

\section{Performance Benchmarks}

\subsection{Target Metrics}

\begin{table}[H]
\centering
\begin{tabularx}{\textwidth}{lXl}
\toprule
\textbf{Metric} & \textbf{Target} & \textbf{Measurement} \\
\midrule
API Response Time (p95) & <200ms & Prometheus \\
WebSocket Message Latency & <50ms & Custom instrumentation \\
Video Stream Latency & <1s & WebRTC stats \\
AI Detection Time & <5s & FastAPI middleware \\
Eye Tracking FPS & 30 FPS & Renderer metrics \\
Database Query Time (p95) & <100ms & pg\_stat\_statements \\
Page Load Time (p75) & <2s & Lighthouse CI \\
\bottomrule
\end{tabularx}
\caption{Performance Targets}
\end{table}

\subsection{Capacity Planning}

\textbf{Baseline Configuration}:
\begin{itemize}[noitemsep]
    \item Supports 1000 concurrent sessions
    \item 100 active interviews simultaneously
    \item 10,000 registered users
    \item 1 TB video storage per month
\end{itemize}

\textbf{Scaling Strategy}:
\begin{itemize}[noitemsep]
    \item \textbf{Horizontal}: HPA scales pods based on CPU (>70\%) and memory (>80\%)
    \item \textbf{Vertical}: Increase pod resources for stateful services
    \item \textbf{Database}: Add read replicas for high query load
    \item \textbf{Redis}: Add nodes to cluster for more memory
    \item \textbf{Storage}: Auto-scaling S3 with lifecycle policies
\end{itemize}

%----------------------------------------------------------------------------------------
% APPENDICES
%----------------------------------------------------------------------------------------

\appendix

\section{API Reference}

\subsection{Authentication Endpoints}

\subsubsection{POST /api/v1/auth/register}

Register a new user account.

\textbf{Request Body}:
\begin{lstlisting}[language=json]
{
  "email": "interviewer@company.com",
  "password": "SecurePass123!",
  "full_name": "Jane Doe",
  "role": "interviewer",
  "organization_id": "uuid"
}
\end{lstlisting}

\textbf{Response} (201 Created):
\begin{lstlisting}[language=json]
{
  "user_id": "uuid",
  "email": "interviewer@company.com",
  "access_token": "eyJhbGc...",
  "refresh_token": "refresh_token_here",
  "expires_in": 3600
}
\end{lstlisting}

\subsubsection{POST /api/v1/auth/login}

Authenticate user and obtain JWT tokens.

\textbf{Request Body}:
\begin{lstlisting}[language=json]
{
  "email": "interviewer@company.com",
  "password": "SecurePass123!",
  "mfa_code": "123456"  // Optional, required if MFA enabled
}
\end{lstlisting}

\textbf{Response} (200 OK):
\begin{lstlisting}[language=json]
{
  "user_id": "uuid",
  "access_token": "eyJhbGc...",
  "refresh_token": "refresh_token_here",
  "expires_in": 3600
}
\end{lstlisting}

\subsection{Session Management Endpoints}

\subsubsection{POST /api/v1/sessions}

Create a new interview session.

\textbf{Request Body}:
\begin{lstlisting}[language=json]
{
  "interviewee_id": "uuid",
  "scheduled_start": "2025-11-25T10:00:00Z",
  "questions": [
    {
      "question_text": "Explain the difference between TCP and UDP",
      "expected_duration_seconds": 180,
      "difficulty": "medium"
    }
  ]
}
\end{lstlisting}

\textbf{Response} (201 Created):
\begin{lstlisting}[language=json]
{
  "session_id": "uuid",
  "session_token": "token_for_interviewee",
  "status": "scheduled",
  "join_url": "https://blocked.com/join/token"
}
\end{lstlisting}

\section{Configuration Reference}

\subsection{Environment Variables}

\begin{longtable}{p{0.3\textwidth}p{0.6\textwidth}}
\toprule
\textbf{Variable} & \textbf{Description} \\
\midrule
\endfirsthead
\toprule
\textbf{Variable} & \textbf{Description} \\
\midrule
\endhead
\texttt{NODE\_ENV} & Environment: development, staging, production \\
\texttt{API\_PORT} & API Gateway port (default: 8000) \\
\texttt{DATABASE\_URL} & PostgreSQL connection string \\
\texttt{REDIS\_URL} & Redis connection string \\
\texttt{RABBITMQ\_URL} & RabbitMQ connection string \\
\texttt{JWT\_PRIVATE\_KEY} & RS256 private key for signing JWTs \\
\texttt{JWT\_PUBLIC\_KEY} & RS256 public key for verifying JWTs \\
\texttt{OPENAI\_API\_KEY} & OpenAI API key for GPT-4 \\
\texttt{ANTHROPIC\_API\_KEY} & Anthropic API key for Claude \\
\texttt{GOOGLE\_AI\_API\_KEY} & Google API key for Gemini \\
\texttt{AWS\_ACCESS\_KEY\_ID} & AWS credentials for S3 \\
\texttt{AWS\_SECRET\_ACCESS\_KEY} & AWS secret key \\
\texttt{S3\_BUCKET\_NAME} & S3 bucket for video storage \\
\bottomrule
\end{longtable}

\section{Troubleshooting}

\subsection{Common Issues}

\subsubsection{API Gateway Returns 502 Bad Gateway}

\textbf{Symptoms}: API requests fail with 502 status code

\textbf{Possible Causes}:
\begin{itemize}
    \item Backend service is down
    \item Database connection pool exhausted
    \item Network policy blocking traffic
\end{itemize}

\textbf{Solutions}:
\begin{lstlisting}[language=bash]
# Check pod status
kubectl get pods -n production

# Check logs
kubectl logs -f deployment/api-gateway -n production

# Check database connections
kubectl exec -it deployment/api-gateway -n production -- \
  psql $DATABASE_URL -c "SELECT count(*) FROM pg_stat_activity;"

# Restart deployment if needed
kubectl rollout restart deployment/api-gateway -n production
\end{lstlisting}

\subsubsection{Eye Tracking Not Working}

\textbf{Symptoms}: No gaze data received from browser

\textbf{Possible Causes}:
\begin{itemize}
    \item Camera permission denied
    \item MediaPipe models not loaded
    \item Calibration not performed
\end{itemize}

\textbf{Solutions}:
\begin{enumerate}
    \item Check browser console for errors
    \item Verify camera permission granted
    \item Run calibration process (9-point grid)
    \item Check MediaPipe model files are present
\end{enumerate}

%----------------------------------------------------------------------------------------
% GLOSSARY
%----------------------------------------------------------------------------------------

\section{Glossary}

\begin{description}
    \item[API Gateway] Centralized entry point for all client requests, handles routing and authentication
    \item[Chromium Fork] Custom version of Chromium browser with embedded security features
    \item[HPA] Horizontal Pod Autoscaler - Kubernetes feature for auto-scaling pods
    \item[JWT] JSON Web Token - Token-based authentication mechanism
    \item[MediaPipe] Google's framework for building multimodal ML pipelines
    \item[mediasoup] WebRTC Selective Forwarding Unit (SFU) for video streaming
    \item[Mojo] Chromium's IPC system for communication between processes
    \item[pgvector] PostgreSQL extension for storing and querying vector embeddings
    \item[SFU] Selective Forwarding Unit - Video routing architecture
    \item[TimescaleDB] PostgreSQL extension optimized for time-series data
    \item[TOTP] Time-based One-Time Password - MFA authentication method
    \item[WebRTC] Web Real-Time Communication - Peer-to-peer communication protocol
    \item[XGBoost] Gradient boosting framework used for AI detection classification
\end{description}

%----------------------------------------------------------------------------------------
% REFERENCES
%----------------------------------------------------------------------------------------

\section{References}

\begin{enumerate}
    \item Chromium Development Documentation: \url{https://www.chromium.org/developers/}
    \item MediaPipe Documentation: \url{https://developers.google.com/mediapipe}
    \item PostgreSQL 18 Documentation: \url{https://www.postgresql.org/docs/18/}
    \item TimescaleDB Documentation: \url{https://docs.timescale.com/}
    \item Kubernetes Documentation: \url{https://kubernetes.io/docs/}
    \item Fastify Documentation: \url{https://www.fastify.io/}
    \item React 19 Documentation: \url{https://react.dev/}
    \item mediasoup Documentation: \url{https://mediasoup.org/documentation/}
\end{enumerate}

%----------------------------------------------------------------------------------------
% VERSION HISTORY
%----------------------------------------------------------------------------------------

\section{Version History}

\begin{table}[H]
\centering
\begin{tabularx}{\textwidth}{llX}
\toprule
\textbf{Version} & \textbf{Date} & \textbf{Changes} \\
\midrule
1.0.0 & 2025-11-24 & Initial release - Complete system documentation \\
\bottomrule
\end{tabularx}
\end{table}

%----------------------------------------------------------------------------------------

\end{document}
